\chapter{通信记录参考}

本附录集中整理通信系统的记录、状态和命令关系。它服务于开发者调试，不替代源码模型定义。

\section{记录实体}

\begin{longtable}{p{0.20\linewidth}p{0.26\linewidth}p{0.44\linewidth}}
\toprule
实体 & 主要来源 & 开发者关注点 \\
\midrule
Submission & message-bureau submission store & 一次 logical submit 的聚合，特别是 broadcast。 \\
Message & \src{MessageRecord} & target scope、target agents、reply policy、retry policy、state。 \\
Attempt & \src{AttemptRecord} & agent、provider、job id、retry index、attempt state。 \\
Job & \src{JobRecord} & dispatcher 执行事实、provider options、terminal decision。 \\
Inbound event & mailbox kernel & task request/reply、head ordering、ack/consume/abandon。 \\
Reply & \src{ReplyRecord} & terminal status、reply text/artifact、diagnostics。 \\
Callback edge & callback edge store & parent、child、continuation、timeout、reply capture。 \\
Completion snapshot & completion storage & tracker state、decision、provider cursor。 \\
\bottomrule
\end{longtable}

\section{常见状态转换}

\subsection{普通 ask}

\begin{Verbatim}[fontsize=\small,frame=single]
Message: created -> running -> completed/failed/incomplete/cancelled
Attempt: queued -> running -> completed/failed/incomplete/cancelled
Job: accepted/queued -> running -> completed/failed/incomplete/cancelled
Inbound task_request: queued -> delivering/claimed -> consumed/abandoned
Reply event: queued -> head -> consumed by ack or reply delivery
\end{Verbatim}

\subsection{retry}

\begin{Verbatim}[fontsize=\small,frame=single]
old attempt: terminal non-completed
same message: append new attempt
same submission: append new job
new job: accepted/queued -> running -> terminal
trace: shows retry_index increasing
\end{Verbatim}

\subsection{resubmit}

\begin{Verbatim}[fontsize=\small,frame=single]
old message: unchanged
new message: origin points to old message
new submission/job chain: created through normal submit plan
trace old and new separately
\end{Verbatim}

\subsection{callback}

\begin{Verbatim}[fontsize=\small,frame=single]
parent job running
child ask submitted with callback mode
callback edge created
child completes or is cancelled
child reply captured
exactly one continuation job submitted to parent agent
edge marked advanced/done
\end{Verbatim}

\section{命令与记录关系}

\begin{longtable}{p{0.20\linewidth}p{0.30\linewidth}p{0.40\linewidth}}
\toprule
命令 & 读取/修改 & 说明 \\
\midrule
\cmd{ask} & 写 submission/message/attempt/job/inbound event & 入口命令。 \\
\cmd{watch} & 读 job events/completion payload & cursor-based event stream。 \\
\cmd{pend} & 读 job get 或 mailbox head overlay & 组合观察。 \\
\cmd{queue} & 读 mailbox summary 和 inbound events & mailbox-centric queue view。 \\
\cmd{inbox} & 读 mailbox head 和 pending events & caller reply view。 \\
\cmd{ack} & 修改 head inbound event 状态 & 消费 reply 或 terminal request。 \\
\cmd{trace} & 读 submission/message/attempt/reply/event/job & lineage graph。 \\
\cmd{retry} & 写新 job 和 retry attempt & 保留原 message。 \\
\cmd{resubmit} & 写新 message/submission/jobs & 新链路。 \\
\cmd{cancel} & 写 cancelled decision 和 terminal attempt & 主动终止。 \\
\bottomrule
\end{longtable}

\section{诊断字段}

trace 的 summary 字段含义：

\begin{itemize}
  \item \src{message_count}：payload 中 message summaries 数量。
  \item \src{attempt_count}：payload 中 attempt summaries 数量。
  \item \src{reply_count}：payload 中 reply summaries 数量。
  \item \src{event_count}：payload 中 inbound event summaries 数量。
  \item \src{job_count}：payload 中 job summaries 数量。
\end{itemize}

queue 的 summary 字段含义：

\begin{itemize}
  \item \src{queue_depth}：mailbox pending queue depth。
  \item \src{pending_reply_count}：pending task reply 数量。
  \item \src{active_inbound_event_id}：当前 mailbox active event。
  \item \src{summary_status}：\src{ok}、\src{missing} 或 \src{error}。
  \item \src{runtime_state} 和 \src{runtime_health}：dispatcher facade 叠加的 registry 视图。
\end{itemize}

inbox 的 head 字段含义：

\begin{itemize}
  \item \src{head_inbound_event_id}：当前可 ack 或 inspect 的 head。
  \item \src{head_event_type}：通常是 task request 或 task reply。
  \item \src{head_reply_id}：如果 head 是 reply，则指向 ReplyRecord。
  \item \src{head_reply_notice}：reply 是否带 notice 元数据。
  \item \src{head_reply_last_progress_at}：reply delivery/progress 相关时间。
\end{itemize}

\section{排障配方}

\subsection{用户说 ask 没有回复}

\begin{enumerate}
  \item 查 \cmd{trace <job_id>}，确认 job 是否 terminal。
  \item 如果 job running，查 \cmd{watch <job_id>} 和 provider runtime。
  \item 如果 job terminal 但无 reply，查 message-bureau finalization。
  \item 如果有 reply 但用户没看到，查 caller \cmd{inbox --detail <caller>}。
  \item 如果 reply head 存在但无法消费，查 reply delivery job 和 ack 限制。
\end{enumerate}

\subsection{queue 显示 degraded}

\begin{enumerate}
  \item 读取 \cmd{queue --detail <agent>} 看 summary status。
  \item 若 \src{missing}，等待 maintenance refresh 或查 mailbox store。
  \item 若 \src{error}，运行 \cmd{doctor storage}。
  \item 用 \cmd{trace} 确认 message/attempt/reply 是否仍完整。
\end{enumerate}

\subsection{callback 卡住}

\begin{enumerate}
  \item 找 child job id 并 trace。
  \item 查 callback edge 是否已有 child reply。
  \item 如果 child terminal 但 parent continuation 未提交，查 finalization callback path。
  \item 如果 parent 已有 outstanding callback，新的 callback 会被拒绝。
\end{enumerate}

\subsection{retry 失败}

\begin{enumerate}
  \item 确认目标是 \src{job_id} 或 \src{attempt_id}。
  \item 确认 attempt 是 terminal。
  \item 确认 attempt 不是 completed。
  \item 确认它是该 agent 最新 attempt。
  \item 确认 target agent runtime 可用。
\end{enumerate}

\section{设计约束}

通信层的主要约束可以简化为五条。

\begin{enumerate}
  \item job 执行事实和 message 通信事实必须分离。
  \item mailbox head ordering 不能被 ack 跳过。
  \item retry 延续 message lineage，resubmit 创建新 lineage。
  \item callback 必须显式建 edge，不能隐式同步等待。
  \item artifact 是 payload indirection，不是新的通信对象类型。
\end{enumerate}

维护者如果守住这五条，通信系统的可解释性就不会轻易丢失。
